\documentclass[aspectratio=169,9pt]{beamer}
\mode<presentation> 
\usepackage[utf8]{inputenc}
\usepackage{sansmathaccent} %Für equations
\usepackage{amsmath}
\usepackage{tikz}
\usepackage{multicol}
\usepackage{cancel}
\usepackage{relsize}
\usepackage{graphicx}
\usepackage[table]{xcolor}
\usepackage{fontawesome5}
\usepackage{enumitem}
\usepackage[labelfont={bf},font={color=black!40,footnotesize},tablename={Tabelle},figurename={Abbildung}]{caption}
\usetikzlibrary{positioning, arrows.meta, calc}
\usefonttheme[onlymath]{serif}
%\makeatletter
%\makeatother

\rowcolors{2}{miugray!30}{miugray!10}

\usetheme{Madrid}  %% Themenwahl
\useoutertheme{split} % Alternatively: miniframes, infolines, split
%\useinnertheme{circles}


\setbeamertemplate{itemize item}{\color{miugray}$\blacksquare$}
\setbeamertemplate{itemize subitem}{\color{uniblau}$\bullet$}



\newcommand{\metermilli}{\mskip3mu \text{mm}}
\newcommand{\cm}{\mskip3mu \text{cm}}
\newcommand{\m}{ \text{m}}
\newcommand{\lite}{ \text{l}}
\newcommand{\kg}{ \text{kg}}
\newcommand{\g}{ \text{g}}
\newcommand{\km}{ \text{km}}
\newcommand{\h}{ \text{h}}
\newcommand{\s}{ \text{s}}
\newcommand{\tonne}{ \text{t}}

\newcommand{\RA}{$\rightarrow$~}
\newcommand{\pp}{\par \vspace*{2mm}}
\newcommand{\ppbig}{\par \vspace*{4mm}}

\definecolor{miudiutex}{HTML}{1d2926}
\definecolor{miugray}{HTML}{b4cccf}
\definecolor{miugray2}{HTML}{4d7365}
\definecolor{white2}{HTML}{f5f7f8}
\definecolor{red}{HTML}{cf1f5c}
\definecolor{green}{HTML}{00A36C}
\definecolor{delphinidin}{HTML}{315794}
\definecolor{uniblau}{HTML}{004e9f}
\definecolor{unigelb}{HTML}{fcba00}
\definecolor{unigrau}{HTML}{909085}
\definecolor{pelargonidin}{HTML}{b33807}
\definecolor{cyanidin}{HTML}{ba0746}
\definecolor{paeonidin}{HTML}{d15ca6}
\definecolor{petunidin}{HTML}{b56fed}
\definecolor{malvidin}{HTML}{8a2d8c}

%\setbeamercolor{palette primary}{bg=pelargonidin!35,fg=white}
%\setbeamercolor{palette secondary}{bg=delphinidin!45,fg=white}
\setbeamercolor{palette primary}{bg=uniblau!70,fg=white}
\setbeamercolor{palette secondary}{bg=uniblau!45,fg=white}
\setbeamercolor{palette tertiary}{bg=unigelb!70,fg=unigrau}
\setbeamercolor{palette quaternary}{bg=miugray,fg=white}
\setbeamercolor{structure}{fg=black} % itemize, enumerate, etc
\setbeamercolor{section in toc}{fg=red} % TOC sections

\setbeamersize{text margin left=1cm,text margin right=1cm}


\useinnertheme{tcolorbox} 
\newcommand{\goodtk}[2]{
\begin{tcolorbox}[
	enhanced,
	colframe=uniblau!75,
	sharp corners,
	boxsep=5pt,
	title=\bfseries \sffamily #1,
	colback=white,
	coltitle=black,
	attach boxed title to top text left={yshift=-\tcboxedtitleheight/2},
	boxed title style={colback=white,colframe=white}
]
	#2
	\end{tcolorbox}
}
	
\newcommand{\antwort}[1]{
\begin{tcolorbox}[
	enhanced,
	colframe=unigrau!65,
	sharp corners,
	boxsep=1pt,
	colback=white,
	coltitle=black,
]
	#1
	\end{tcolorbox}
}


\title{3. Übung}
\author{PD Dr. Helene Loos}
\date{\today}

\begin{document}
%\maketitle
\small



\begin{frame}
\begin{center}

\LARGE Übung 4  \par \normalsize
Prozessbezogene Lebensmitteltechnologie
\end{center}
\small
\end{frame}










\begin{frame}




\begin{minipage}{.4\textwidth}
{\color{miugray}$\blacksquare ~$} \textbf{Aufgabe 1} \par \vspace*{2mm}
Berechnen Sie die Strahlungsleistung, die von 100 m$^2$ einer polierten Eisenoberfläche
(Emissionsgrad = 0,06) abgegeben wird. Die Temperatur der Oberfläche beträgt 37 °C.


\end{minipage}%
\hspace*{4mm}
\begin{minipage}{.6\textwidth}
\antwort{
\footnotesize
\begin{align*}
\uncover<2->{P &= \varepsilon \cdot \sigma \cdot A \cdot T^4 \\}
\uncover<3->{P &= 0,06 \cdot  100~\text{m}^2 \cdot 5,67 \cdot 10^{-8} ~\text{W} \cdot \text{m}^{-2}\text{K}^{-4}\cdot 310,15~\text{K}^4 \\}
\uncover<4->{ &= 3.147,9~\text{W}}
\end{align*}
}
\end{minipage}%

\end{frame}

\begin{frame}
\begin{minipage}{.5\textwidth}
{\color{miugray}$\blacksquare ~$} \textbf{Aufgabe 2} \par \vspace*{2mm}
Eine Seite einer 1 cm dicken Edelstahlplatte wird auf 110 °C gehalten, die andere Seite
auf 90 °C. Berechnen Sie unter der Annahme stationärer Bedingungen die
Wärmeübertragungsrate pro Flächeneinheit durch die Platte. Die Wärmeleitfähigkeit von
nichtrostendem Stahl beträgt \mbox{17 W m$^{-1}$ °C$^{-1}$}.

\end{minipage}%
\hspace*{4mm}
\begin{minipage}{.5\textwidth}
\antwort{
\begin{align*}
\uncover<2->{\dot{Q} &= - A \cdot \frac{\lambda}{d} \cdot (T_2 - T_1) \\}
\uncover<3->{- \frac{\dot{Q}}{A} &= \cdot \frac{\lambda}{d} \cdot (T_2 - T_1) \\}
\uncover<4->{- \frac{\dot{Q}}{A} &= \cdot \frac{17 \text{W} / \text{m°C}}{0,01~\text{m}} \cdot 20\text{°C}\\
	&= 34.000 ~\text{W/m$^2$}}
\end{align*}

}
\end{minipage}%
\end{frame}


\begin{frame}

\begin{minipage}{.5\textwidth}
{\color{miugray}$\blacksquare ~$} \textbf{Aufgabe 3} \par \vspace*{2mm}
Ein Glühofen hat eine 4 m$^2$ große Wand aus Schamotte. Die Dicke der Wand sei 0,5 m.
Zwischen der Innen- und der Außenseite der Wand besteht eine Temperaturdifferenz von
1.000 K. Die Temperatur außerhalb des Ofens beträgt 20 °C. Die Wärmeleitfähigkeit der
Schamotte ist $\lambda_1$ = 0,93 W/mK. 

%\begin{figure}
%\centering
%\begin{tikzpicture}
%\draw[fill=miugray!50] 	(0,0) rectangle (0.8,3);
%\draw[fill=white]		(0.8,0) rectangle (2,3);
%\node at (0.4,1.5) {$\dot{Q_1}$};
%\node at (1.4,1.5) {$\dot{Q_2}$};
%\draw[->,shorten >=2pt] 	(0,-0.3) -- (0,0)	node[below,yshift=-2mm]{$T_1$}; 
%\draw[->,shorten >=2pt] 	(0.8,-0.3) -- (0.8,0)node[below,yshift=-2mm]{$T_2$};
%\draw[->,shorten >=2pt] 	(2,-0.3) -- (2,0)	node[below,yshift=-2mm]{$T_3$};
%\end{tikzpicture}
%\caption{Der Wärmestrom ist an jeder Stelle der Wand konstant ($\dot{Q_1} = \dot{Q_2}$). Die Differenz der Temperatur $\Delta T$ lässt sich in ihre Bestandteile $T_1$, $T_2$ und $T_3$ zerlegen.}
%\end{figure}
\pp
\begin{enumerate}[label=\alph*)]

\item Wie groß ist der tägliche Wärmeverlust durch die Wand?

\end{enumerate}

\end{minipage}%
\hspace*{4mm}
\begin{minipage}{.5\textwidth}



\antwort{

\begin{align*}
\uncover<2->{\dot{Q} &= - A \cdot \frac{\lambda}{d} \cdot (T_2 - T_1) \\}
\uncover<3->{&= - 4~\text{m}^2 \cdot \frac{0,93}{0,5~\text{m}} \cdot 1.000 K \\
&= - 7.440~\text{W} \\}
\uncover<4->{&= - 7.440~\text{W} \cdot 24~\text{h} \\
&= - 178.560~\text{Wh}}
\end{align*}
}


\end{minipage}%
\end{frame}



\begin{frame}

\begin{minipage}{.5\textwidth}
{\color{miugray}$\blacksquare ~$} \textbf{Aufgabe 3} \par \vspace*{2mm}
Ein Glühofen hat eine 4 m$^2$ große Wand aus Schamotte. Die Dicke der Wand sei 0,5 m.
Zwischen der Innen- und der Außenseite der Wand besteht eine Temperaturdifferenz von
1.000 K. Die Temperatur außerhalb des Ofens beträgt 20 °C. Die Wärmeleitfähigkeit der
Schamotte ist $\lambda_1$ = 0,93 W/mK. 

\pp 
\begin{enumerate}[label=\alph*)]

%\item Wie groß ist der tägliche Wärmeverlust durch die Wand?
\item[b)] Um wie viel Prozent verringert sich dieser Verlust, wenn die Außenwand (von außen)
durch eine 3 cm dicke Isolierschicht mit einer Wärmeleitfähigkeit von $\lambda_2$ = 0,093 
W/mK verstärkt wird?
%\item Welche Temperatur herrscht an der Grenze zwischen der Isolierung und der
%Schamotte-Schicht?
\end{enumerate}

\end{minipage}%
\hspace*{4mm}
\begin{minipage}{.5\textwidth}



%\antwort{

%\begin{align*}
%\dot{Q} &= - A \cdot \frac{\lambda}{d} \cdot (T_2 - T_1) \\
%\uncover<3>{&= - 4m^2 \cdot \frac{0,93}{0,5~\text{m}} \cdot 1000 K \\
%&= - 7.440~W \\
%&= - 7.440~W \cdot 24h \\
%&= - 178.560~Wh}
%\end{align*}
%}



%
%
\antwort{
\uncover<2->{
\begin{align*}
\dot{Q} &= - \frac{A \cdot \Delta T_{ges}}{(\frac{d_2}{\lambda_2}+\frac{d_1}{\lambda_1})} \\
\end{align*}
}
\uncover<3->{
Einsetzen:
}
\uncover<4->{
\begin{align*}
\dot{Q} &= - \frac{4~\text{m}^2 \cdot 1.000 K}{(\frac{0,03~\text{m}}{0,093 ~\text{W/mk}}+\frac{0,5~\text{m}}{0,93 ~\text{W/mk}})} \\
&= 4.650 ~\text{W}
\end{align*}
}
\uncover<5->{
Prozentualer Unterschied:
}
\uncover<6->{
\begin{align*}
\frac{4.650 ~\text{W}}{7.440 ~\text{W}} &\approx  0,625 = 62,5\%
\end{align*}}
\uncover<7->{
Also verringert sich der Verlust um 100-62,5\% = 37,5\%
}

}
%%
%
%
%
%
%
%
%
%\antwort{
%\begin{align*}
%\dot{Q} &= - \lambda \cdot A \cdot \frac{T_3 - T_2}{d} \\
%T_2 &= \frac{\dot{Q} \cdot d}{\lambda \cdot A } + T_3 \\
%\\
%T_2 &= \frac{4.650 ~\text{W} \cdot 0,03~\text{m}}{0,093 ~\text{W/mk} \cdot 4~\text{m}^2} + 20~^{\circ}\text{C} \\
% &= 395~^{\circ}\text{C}
%\end{align*}
%}
%

\end{minipage}%
\end{frame}


\begin{frame}

\begin{minipage}{.5\textwidth}
{\color{miugray}$\blacksquare ~$} \textbf{Aufgabe 3} \par \vspace*{2mm}
Ein Glühofen hat eine 4 m$^2$ große Wand aus Schamotte. Die Dicke der Wand sei 0,5 m.
Zwischen der Innen- und der Außenseite der Wand besteht eine Temperaturdifferenz von
1.000 K. Die Temperatur außerhalb des Ofens beträgt 20 °C. Die Wärmeleitfähigkeit der
Schamotte ist $\lambda_1$ = 0,93 W/mK. 

%\begin{figure}
%\centering
%\begin{tikzpicture}
%\draw[fill=miugray!50] 	(0,0) rectangle (0.8,3);
%\draw[fill=white]		(0.8,0) rectangle (2,3);
%\node at (0.4,1.5) {$\dot{Q_1}$};
%\node at (1.4,1.5) {$\dot{Q_2}$};
%\draw[->,shorten >=2pt] 	(0,-0.3) -- (0,0)	node[below,yshift=-2mm]{$T_1$}; 
%\draw[->,shorten >=2pt] 	(0.8,-0.3) -- (0.8,0)node[below,yshift=-2mm]{$T_2$};
%\draw[->,shorten >=2pt] 	(2,-0.3) -- (2,0)	node[below,yshift=-2mm]{$T_3$};
%\end{tikzpicture}
%\caption{Der Wärmestrom ist an jeder Stelle der Wand konstant ($\dot{Q_1} = \dot{Q_2}$). Die Differenz der Temperatur $\Delta T$ lässt sich in ihre Bestandteile $T_1$, $T_2$ und $T_3$ zerlegen.}
%\end{figure}
\pp
\begin{enumerate}[label=\alph*)]

%\item Wie groß ist der tägliche Wärmeverlust durch die Wand?
%\item Um wie viel Prozent verringert sich dieser Verlust, wenn die Außenwand (von außen)
%durch eine 3 cm dicke Isolierschicht mit einer Wärmeleitfähigkeit von $\lambda_2$ = 0,093 
%W/mK verstärkt wird?
\item[c)] Welche Temperatur herrscht an der Grenze zwischen der Isolierung und der
Schamotte-Schicht?
\end{enumerate}

\end{minipage}%
\hspace*{4mm}
\begin{minipage}{.5\textwidth}


%
%\antwort{
%
%\begin{align*}
%\dot{Q} &= - A \cdot \frac{\lambda}{d} \cdot (T_2 - T_1) \\
%\uncover<3>{&= - 4m^2 \cdot \frac{0,93}{0,5~\text{m}} \cdot 1000 K \\
%&= - 7.440~W \\
%&= - 7.440~W \cdot 24h \\
%&= - 178.560~Wh}
%\end{align*}
%}



%
%
%\antwort{
%\begin{align*}
%\dot{Q} &= - \frac{A \cdot \Delta T_{ges}}{(\frac{d_2}{\lambda_2}+\frac{d_1}{\lambda_1})} \\
%\end{align*}
%Einsetzen:
%\begin{align*}
%\dot{Q} &= - \frac{4m^2 \cdot 1000 K}{(\frac{0,03m}{0,093 W/mk}+\frac{0,5m}{0,93 W/mk})} \\
%&= 4.650 W
%\end{align*}
%Prozentualer Unterschied:
%\begin{align*}
%\frac{4.650 W}{7.440 W} &\approx  0,625 = 62,5\%
%\end{align*}
%Also verringert sich der Verlust um 100-62,5\% = 37,5\%
%
%
%}
%%







\antwort{
\begin{align*}
\uncover<2->{\dot{Q} &= - \lambda \cdot A \cdot \frac{T_3 - T_2}{d} \\}
\uncover<3->{T_2 &= \frac{\dot{Q} \cdot d}{\lambda \cdot A } + T_3 \\}
\\
\uncover<4->{T_2 &= \frac{4.650 ~\text{W} \cdot 0,03~\text{m}}{0,093 ~\text{W/mk} \cdot 4~\text{m}^2} + 20~^{\circ}\text{C} \\}
\uncover<5->{ &= 395~^{\circ}\text{C}}
\end{align*}
}


\end{minipage}%
\end{frame}




\begin{frame}
\begin{minipage}{.5\textwidth}
{\color{miugray}$\blacksquare ~$} \textbf{Aufgabe 4} \par \vspace*{2mm}
Für die Berechnung der spezifischen Wärmekapazität $c$ kann folgende Formel verwendet
werden:

\begin{align*}
c = &\quad ~   4,180 \times \text{Wasser} \\& + 1,711 \times \text{Proteine} \\& + 1,928 \times \text{Fett} \\&+ 1,547 \times \text{Kohlenhydrate} \\& + 0,908 \times \text{weitere Bestandteile}
\end{align*}


Welche Wärmemenge (in kWh) benötigen Sie, um 6,5 L einer Kartoffelsuppe von 10 °C auf
60 °C zu erhitzen? Es wird angenommen, dass die Dichte der Suppe 1 kg/dm$^3$ beträgt und
es keine Wärmeverluste gibt. \pp 

Die Suppe setzt sich pro 100 g zusammen aus: \par 
\begin{itemize}
\item Fett: 14,3 g
\item Proteine: 2,7 g
\item Kohlenhydrate: 11,5 g
\item Weitere Bestandteile: 0,5 g
\end{itemize}

\end{minipage}%
\hspace*{4mm}
\begin{minipage}{.5\textwidth}
\antwort{
\uncover<2->{
Zunächst Wärmekapazität $c$ ausrechnen:}
\uncover<3->{
\begin{align*}
c &= ~~~ 4,180 \cdot 0,71 + 1,711 \cdot 0,027 + 1,928 \cdot 0,143 \\
&~~~ + 1,547 \cdot 0,115 + 0,908 \cdot 0,005 \\
 &= ~~~ 3,472146~ \text{kJ} \cdot \text{kg}^{-1} \cdot \text{K}^{-1}
\end{align*}
}
\uncover<4->{
Energie berechnen:}

\begin{align*}
\uncover<5->{Q &= m \cdot c \cdot \Delta T \\ \\}
\uncover<6->{Q &= 6,5~ \text{kg} \cdot 3,472~ \text{kJ}/\text{kg K} \cdot (333,15-283,15 ~ \text{K}) \\
&= 1.128,45~ \text{kJ} \\
&= 1.128,45~ \text{kWs}  \\
&= 0,313~ \text{kWh}}
\end{align*}

}
\end{minipage}%
\end{frame}






\begin{frame}
\begin{minipage}{.5\textwidth}
{\color{miugray}$\blacksquare ~$} \textbf{Aufgabe 5} \par \vspace*{2mm}
In Ihrem 20 °C warmen Zuhause schalten Sie Ihre Herdplatte auf maximaler Stufe an. Diese erhitzt sich daraufhin auf 300 °C. Im nächsten Schritt stellen Sie eine Pfanne aus Aluminium ($\lambda_{Al} = 236~\text{W/mK}$) mit einem Durchmesser von 28 cm und einer Dicke von 2 mm auf die Herdplatte.
\end{minipage}%
\hspace*{4mm}
\begin{minipage}{.5\textwidth}
\begin{itemize}
\item Welcher Wärmefluss entsteht zwischen der Unter- und der Oberseite der Pfanne?

\antwort{
\begin{align*}
\uncover<2->{\dot{Q} &= \lambda_{Al} \cdot \frac{A}{d} \cdot (\Delta T) \\}
\uncover<3->{
&= 236~\text{W/mK} \cdot \frac{\pi \cdot 0,14~\text{m}^2}{0,002~\text{m}} \cdot 280~\text{K}\\
&= 2.034.445~\text{J}\\
&= 2.034~\text{kJ}}
\end{align*}
}

\item Wieso können Sie von dem fließenden Wärmestrom keine Rückschlüsse auf ihre tatsächliche Stromrechnung schließen?

\antwort{
\uncover<4->{
Das Delta T wird mit der Zeit immer geringer. Die Herdplatte muss anfänglich eine viel größere Leistung aufbringen um die kalte Herdplatte (bzw auch die Pfanne) zu erhitzen. Ein typischer Herd hat eine Leistung von circa 10 kW. Für einen konstanten Wärmestrom aus unserer Aufgabe wären 2 Megawatt notwendig!
}}

\end{itemize}
\end{minipage}%
\end{frame}



\begin{frame}
\begin{minipage}{.4\textwidth}

{\color{miugray}$\blacksquare ~$} \textbf{Aufgabe 6} \par \vspace*{2mm}
Eine gefrorene Kartoffelsuppe (-19 °C, 850 g) soll mithilfe einer Mikrowelle zum Verzehr auf
80 °C erwärmt werden. Die Leistung der Mikrowelle liegt bei 800 W. Nehmen Sie an, dass
diese Leistung zu 100\% für die Erwärmung der Kartoffelsuppe zur Verfügung steht.
Die spezifische Wärmekapazität der gefrorenen Suppe beträgt 2,0 kJ/(kgK) und die der
aufgetauten Suppe beträgt 4,62 kJ/(kgK). Es werden eine spezifische Schmelzwärme von 334
kJ/kg und eine Dichte der Suppe von 1 g/cm$^3$ angenommen. \pp
Wie lange dauert es, bis die gewünschte Temperatur erreicht ist?

\end{minipage}%
\hspace*{4mm}
\begin{minipage}{.6\textwidth}
\antwort{\footnotesize
\begin{align*}
\uncover<2->{Q &= m \cdot c \cdot \Delta T \\}
\uncover<3->{Q &= 0,85~\text{kg} \cdot 2 ~\text{kJ/kgK} \cdot 19 ~\text{K} + 0,85~\text{kg} \cdot 334 ~\text{kJ/kg} 
\\&~~~ + 0,85~\text{kg} \cdot 4,62 ~\text{kJ/kgK} \cdot 80 ~\text{K}\\
	&= 630,36 ~\text{kJ}  \\}
\uncover<4->{	&= 630.360 ~\text{Ws }}\uncover<5->{ && | : 800 ~\text{W} \\}
\uncover<6->{	&\approx 787~\text{s} \\
	&\approx 13~\text{min}~ 8~\text{s}}
\end{align*}

}
\end{minipage}%
\end{frame}







\begin{frame}
\begin{minipage}{.5\textwidth}
{\color{miugray}$\blacksquare ~$} \textbf{Aufgabe 7} \par \vspace*{2mm}
Auf welche Temperatur kühlt sich ein Getränk (500 g, 20 °C) ab, wenn 50 g Eis mit einer
Temperatur von -20 °C hinzugegeben werden? Die spezifische Wärmekapazität von Eis liegt
bei 2,0 kJ/(kgK) und von Wasser bei 4,2 kJ/(kg K). Die spezifische Schmelzwärme von Eis ist angegeben mit $q_S$ = 334 kJ/kg. Vernachlässigen Sie etwaige Wärmeströme in Bezug auf das
Glas oder die Luft. \pp
Berechnen Sie die Endtemperatur $T_f$. Überlegen Sie dafür zunächst, in welche (berechenbare) Schritte Sie diesen Prozess aufteilen können.
\end{minipage}%
\hspace*{4mm}
\begin{minipage}{.5\textwidth}
\antwort{
\uncover<2->{
Zunächst Energie des Eiswürfels ausrechnen (wir gehen davon aus, dass er schmelzen wird)
}
\uncover<3->{
\begin{align*}
Q &= m \cdot c \cdot \Delta T \\ 
Q &= 0,05 ~\text{kg} \cdot 2 ~\text{kJ/kg} \cdot 20 ~\text{K} + 0,05 \cdot 334 ~\text{kJ/kg} \\ 
Q &= 18,7 ~\text{kJ} 
\end{align*}
}
\uncover<4->{
Im nächsten Schritt schauen wir uns an, wie sehr wie die Temperatur unseres Getränkes mit der Energie aus dem Eiswürfel absenken können:
}

\begin{align*}
\uncover<5->{Q &= m \cdot c \cdot \Delta T}\uncover<6->{ \qquad \rightarrow \qquad \Delta T = \frac{Q}{m \cdot c} \\ \\}
\uncover<7->{\Delta T &= \frac{18,7 ~\text{kJ}}{0,5 ~\text{kg} \cdot 4,2 ~\text{kJ/kgK}} \\ 
		&= 8,9 ~\text{°C}}
\end{align*}

}



\end{minipage}%
\end{frame}


\begin{frame}
\begin{minipage}{.5\textwidth}
{\color{miugray}$\blacksquare ~$} \textbf{Aufgabe 7} \par \vspace*{2mm}
Auf welche Temperatur kühlt sich ein Getränk (500 g, 20 °C) ab, wenn 50 g Eis mit einer
Temperatur von -20 °C hinzugegeben werden? Die spezifische Wärmekapazität von Eis liegt
bei 2,0 kJ/(kgK) und von Wasser bei 4,2 kJ/(kg K). Die spezifische Schmelzwärme von Eis ist angegeben mit $q_S$ = 334 kJ/kg. Vernachlässigen Sie etwaige Wärmeströme in Bezug auf das
Glas oder die Luft. \pp
Berechnen Sie die Endtemperatur $T_f$. Überlegen Sie dafür zunächst, in welche (berechenbare) Schritte Sie diesen Prozess aufteilen können.
\end{minipage}%
\hspace*{4mm}
\begin{minipage}{.5\textwidth}
\antwort{

\uncover<2->{
Wir wissen nun, dass unser Getränk eine Temperatur von 11,1°C (20-8,9 °C) und unsere Flüssigkeit, die ehemals Eiswürfel war, eine Temperatur von 0°C hat. Wenn zwei Flüssigkeiten mit unterschiedlichen Temperaturen gemischt werden, gilt für die Endtemperatur: 
}

\begin{align*}
\uncover<2->{T_f &= \frac{c_1 \cdot m_1 \cdot T_1 + c_2 \cdot m_2 \cdot T_2}{c_1 \cdot m_1 + c_2 \cdot m_2} \\}
\uncover<3->{T_f &= \frac{m_1 \cdot T_1 + m_2 \cdot T_2}{m_1 + m_2} \\}
\uncover<4->{T_f &= \frac{0,5~\text{kg} \cdot 11,1 ~\text{°C} + 0,05~\text{kg} \cdot 0 ~\text{°C}}{0,5 ~\text{kg} + 0,05 ~\text{kg}} \\}
\uncover<5->{	&= 10,09 ~\text{°C}}
\end{align*}

}



\end{minipage}%
\end{frame}


\begin{frame}
\begin{center}
Ende Teil 1
\end{center}
\end{frame}



\end{document}
